% Generated draft: requires mathematical and source review.
\documentclass[11pt]{article}
\usepackage[margin=25mm]{geometry}
\usepackage{fontspec}
\setmainfont{Latin Modern Roman}
\usepackage{amsmath,amssymb,mathtools,unicode-math}
\setmathfont{Latin Modern Math}
\usepackage{enumitem,needspace,xurl,hyperref}
\setlength{\parindent}{0pt}
\setlength{\parskip}{5pt}
\setlength{\emergencystretch}{3em}
\newcommand{\N}{\mathbb{N}}
\newcommand{\Z}{\mathbb{Z}}
\newcommand{\Q}{\mathbb{Q}}
\newcommand{\R}{\mathbb{R}}
\newcommand{\C}{\mathbb{C}}
\newcommand{\F}{\mathbb{F}}
\newcommand{\A}{\mathbb{A}}
\newcommand{\PP}{\mathbb{P}}
\newcommand{\E}{\mathbb{E}}
\newcommand{\eps}{\varepsilon}
\newcommand{\dd}{\,\mathrm{d}}
\begin{document}
\section*{UnsolvedMath ID: 30004609}
\textbf{Source problem code:} catalogue entry
\subsection{Definitions and mathematical statement}
Let \(K\) denote the intended base field, which the supplied record does not specify. Let \(V\) be a finite-dimensional \(K\)-vector space of dimension \(\ell\). A central hyperplane arrangement \(\mathcal A\) in \(V\) is a finite set of linear hyperplanes. Write \(S=\operatorname{Sym}(V^*)\), and for each \(H\in\mathcal A\), choose a nonzero linear form \(\alpha_H\) with \(H=\ker\alpha_H\). Its module of logarithmic derivations is
\[
D(\mathcal A)=\{\theta\in\operatorname{Der}_K(S):\theta(\alpha_H)\in\alpha_H S\text{ for every }H\in\mathcal A\}.
\]
The arrangement is free if \(D(\mathcal A)\) is a free \(S\)-module. For a homogeneous basis \(\theta_1,\ldots,\theta_\ell\), its exponents are the multiset \(\exp(\mathcal A)=\{d_1,\ldots,d_\ell\}\), where \(d_i\) is the polynomial degree of the coefficients of \(\theta_i\).

Let \(L(\mathcal A)\) be the lattice of intersections of subfamilies of \(\mathcal A\), including the empty intersection \(V\), ordered by reverse inclusion. Set \(\rho(X)=\operatorname{codim}_V X\). An element \(X\in L(\mathcal A)\) is modular if
\[
\rho(X)+\rho(Y)=\rho(X\wedge Y)+\rho(X\vee Y)
\qquad\text{for every }Y\in L(\mathcal A),
\]
where meet and join are taken in this lattice. Put \(T=\bigcap_{H\in\mathcal A}H\) and \(r=\operatorname{codim}_V T\). The arrangement is supersolvable if there is a chain
\[
V=X_0\supsetneq X_1\supsetneq\cdots\supsetneq X_r=T
\]
of modular elements of \(L(\mathcal A)\) with \(\rho(X_i)=i\).

The supplied conjecture asks whether, for every finite-dimensional \(V\) and every central free arrangement \(\mathcal A\) over the intended field \(K\),
\[
\left(d_i\in\{1,2,3\}\text{ for every }i,\quad
\#\{i:d_i=3\}\leq 1\right)
\quad\Longrightarrow\quad
\mathcal A\text{ is supersolvable}.
\]
There is no rank bound in the supplied assertion. The intended base field, its characteristic, and whether the assertion quantifies over fields remain unresolved; the displayed formulation does not settle those missing qualifications. The exponents here are those of the ambient derivation module, with no zero exponents permitted by the stated hypothesis; the source does not explicitly discuss essentialization.
\subsection{Short English statement}
Must every central free hyperplane arrangement with all exponents equal to \(1\) or \(2\), except for at most one exponent equal to \(3\), have a maximal chain of modular elements in its intersection lattice? The supplied conjecture has no rank bound, but its intended base field and characteristic require clarification.
\subsection{Sources}
\begin{itemize}
\item Supplied local UnsolvedMath registry record, permanent ID \(30004609\), source problem code \texttt{catalogue entry}, titled ``Supersolvability of Low-Exponent Free Hyperplane Arrangements.'' The record belongs to the Oberwolfach Reports Open Problems collection, gives proposed year 2021, and supplies no individual proposer.
\item Supplied source citation: \emph{Logarithmic Vector Fields and Freeness of Divisors and Arrangements: New perspectives and applications} (2022), Oberwolfach Reports, pp.\ 225--300, DOI \texttt{10.4171/owr/2021/5}. \url{https://doi.org/10.4171/owr/2021/5}. The record describes these reports as documenting mathematical workshops and their open-problem discussions.
\item Supplied literature reference: Tohaneanu, \emph{Free arrangements with low exponents}. \url{https://arxiv.org/abs/1707.07091}.
\item Supplied literature reference: \emph{Oberwolfach Report 5/2021}. \url{https://ems.press/content/serial-article-files/46883}.
\item Literature assessment dated \texttt{2026-08-22}: the supplied curation labels the problem open and states that the conjecture is proved in several ranks and for some exponent patterns, while the arbitrary-rank assertion remains open. This is supplied dated literature triage derived from the OWR curation, not an independently verified assessment. Its time-sensitive status and cited claims require independent verification.
\item The registry labels the statement \texttt{corrected\_verified} and reports a manual semantic, typographic, and LaTeX review confirming preservation of the mathematical target. These are supplied editorial qualifications, not independent verification of this draft or of the conjecture's mathematical status.
\end{itemize}
\end{document}
